\textit{Objective.} EEG classification of Alzheimer's disease (AD) and frontotemporal dementia (FTD) on one
public dataset reports accuracies from about 50\% to 99\%, but studies change representation, graph, model and
evaluation protocol at once. We isolated representation and graph construction under a fixed graph
neural network, and tested a candidate EEG marker linked to blood--brain-barrier dysfunction.
\textit{Approach.} On OpenNeuro ds004504 (36 AD, 23 FTD, 29 controls) we crossed seven node-feature
representations with three graph constructions under one compact, untuned graph convolutional network, on 25
frozen subject-level splits (\nMainRuns{} runs) and 25 confirmation splits, with a by-subject mixed
model, overlap-corrected and equivalence tests, matched non-graph and deep-learning baselines, 16 ablations,
permutation and leakage experiments, and external validation on a Parkinson's disease cohort (ds004584,
$n=149$). We also characterised paroxysmal slow-wave events (PSWEs).
\textit{Main results.} Representation had the largest effect in every analysis (by-subject mixed model
$F(\mixSelRepDf,\mixSelRepDenDf)=\mixSelRepF$; \mixConfRepF{} on confirmation splits; \mixExtRepF{} externally)
and its ordering replicated across independent splits ($\rho=\confRankRho$); allowing for overlapping training
sets, it remained significant externally ($p\nbExtRepP$) but not on ds004504 ($p\nbSelRepP$; confirmation $p\nbConfRepP$). Graph construction
had no detectable effect, averaged graph types were equivalent within $\pm\eqMargin$ macro-F1, and a learned
graph did no better. The best configuration (wavelet features, spatial graph) reached macro-F1 \confFOne{}
(95\% CI \confFOneCiLo{}--\confFOneCiHi{}; permutation $p=\cpermPThreespatialP$); non-graph and deep-learning
baselines did not differ significantly. Window-level splitting raised macro-F1 to
\cleakPThreespatialLeakySubj{}. The representation ranking did not transfer to Parkinson's disease
($\rho=\extRepsRho$). PSWEs were more frequent in AD (rate ratio \psweAdjADRR{}) and FTD (\psweAdjFTDRR{}) but
not in Parkinson's disease, and the excess tracked background slowing ($\rho=\psweSlowRho$) under every event
rule tested; among patients whose slowing overlapped with controls' it was not detectable.
\textit{Significance.} The evaluation protocol moves reported performance far more than any modelling choice
we tested, representation matters more than graph construction, and the best representation is
disease-specific. Fixed-threshold PSWEs on scalp EEG mainly index background slowing. Code, frozen folds and
all raw runs are released.
